\subsection{APPLICATIONS: II. MULTIPLICATION OF RATIONAL INTEGERS}

Lemma 2. Let \(\mathbf{E}\) be a monoid and \(x\) an element of \(\mathbf{E}\) .

\begin{enumerate}
\def\labelenumi{(\roman{enumi})}
\item
  There exists a unique homomorphism \(f\) of \(\mathbf{N}\) into \(\mathbf{E}\) with \(f(1) = x\) and \(f(n) = \frac{n}{T} x\) for all \(n \in \mathbf{N}\) .
\item
  If \(x\) is invertible, there exists a unique homomorphism \(g\) of \(\mathbf{Z}\) into \(\mathbf{E}\) such that \(g(1) = x\) and \(g\) coincides with \(f\) on \(\mathbf{N}\) .
\end{enumerate}

Writing \(f(n) = \top^n x\) for all \(n \in \mathbf{N}\) , the formulae

\[
\overset {0} {\top} x = e \quad \text { and } \quad \left(\overset {m} {\top} x\right) \top \left(\overset {n} {\top} x\right) = \overset {m + n} {\top} x
\]

(no. 1) express the fact that \(f\) is a homomorphism of \(\mathbf{N}\) into \(\mathbf{E}\) and obviously \(f(1) = x\) . If \(f'\) is a homomorphism of \(\mathbf{N}\) into \(\mathbf{E}\) such that \(f'(1) = x\) , then \(f = f'\) , by § 1, no. 4, Proposition 1, (iv).

Suppose now that \(x\) is invertible. By no. 3, Corollary 2 to Proposition 4, \(f(n) = \overset{n}{\top} x\) is invertible for every integer \(n \geqslant 0\) . By construction, \(\mathbf{Z}\) is the group of differences of \(\mathbf{N}\) and hence (no. 4, Theorem 1) \(f\) extends uniquely to a homomorphism \(g\) of \(\mathbf{Z}\) into \(\mathbf{E}\) . If \(g'\) is a homomorphism of \(\mathbf{Z}\) into \(\mathbf{E}\) with \(g'(1) = x\) , the restriction \(f'\) of \(g'\) to \(\mathbf{N}\) is a homomorphism of \(\mathbf{N}\) into \(\mathbf{E}\) with \(f'(1) = x\) . Hence \(f' = f\) , whence \(g' = g\) .

We shall apply Lemma 2 to the case where the monoid E is Z; for every integer \(m \in Z\) there therefore exists an endomorphism \(f_{m}\) of Z characterized by \(f_{m}(1) = m\) . If m is in N, the mapping \(n \mapsto mn\) of N into N is an endomorphism of the magma N (Set Theory, III, §3, no. 3, Corollary to Proposition 5); hence \(f_{m}(n) = mn\) for all m, n in N.

Multiplication on \(\mathbf{N}\) can therefore be extended to multiplication on \(\mathbf{Z}\) by the formula \(mn = f_m(n)\) for \(m, n\) in \(\mathbf{Z}\) . We shall establish the formulae:

\[
x y = y x\tag{2}
\]

\[
(x y) z = x (y z)\tag{3}
\]

\[
x (y + z) = x y + x z\tag{4}
\]

\[
(x + y) z = x z + y z\tag{5}
\]

\[
0. x = x. 0 = 0\tag{6}
\]

\[
1. x = x. 1 = x\tag{7}
\]

\[
(- 1). x = x. (- 1) = - x\tag{8}
\]

for \(x, y, z\) in \(\mathbf{Z}\) . (*In other words, \(\mathbf{Z}\) is a commutative ring.*) The formulae \(x(y + z) = xy + xz\) and \(x.0 = 0\) express the fact that \(f_x\) is an endomorphism of the additive monoid \(\mathbf{Z}\) and \(f_{x}(1) = x\) may be written \(x.1 = x\) . The endomorphism \(f_{x} \circ f_{y}\) of \(\mathbf{Z}\) maps 1 to \(xy\) and hence is equal to \(f_{xy}\) , whence (3). Now \(f_{x}(-y) = -f_{x}(y)\) , that is \(x(-y) = -xy\) ; similarly, the endomorphism \(y \mapsto -xy\) of \(\mathbf{Z}\) maps 1 to \(-x\) , whence \((-x).y = -xy\) and therefore

\[
(- x) (- y) = - (x (- y)) = - (- x y) = x y.
\]

For \(m, n\) in \(\mathbf{N}\) , \(mn = nm\) (Set Theory, III, § 3, no. 3, Corollary to Proposition 5), whence \((-m).n = n(-m)\) and \((-m)(-n) = (-n)(-m)\) ; as \(\mathbf{Z} = \mathbf{N} \cup (-\mathbf{N})\) , \(xy = yx\) for \(x, y\) in \(\mathbf{Z}\) ; and this formula means that (5) follows from (4) and completes the proof of formula (6) to (8).

